\documentclass[conference]{IEEEtran}
\IEEEoverridecommandlockouts

% --------------------
% Packages
% --------------------
\usepackage{cite}
\usepackage{amsmath,amssymb,amsfonts}
\usepackage{algorithm}
\usepackage{algorithmic}
\usepackage{graphicx}
\usepackage{xcolor}
\usepackage{booktabs}
\usepackage{textcomp}

% --------------------
% Notation helpers
% --------------------
\DeclareMathOperator*{\argmin}{arg\,min}
\DeclareMathOperator*{\argmax}{arg\,max}
\newcommand{\E}{\mathbb{E}}
\newcommand{\Prob}{\mathbb{P}}
\newcommand{\norm}[1]{\left\lVert #1 \right\rVert}
\newcommand{\abs}[1]{\left| #1 \right|}
\newcommand{\set}[1]{\left\{ #1 \right\}}
\newcommand{\given}{\,\mid\,}
\newcommand{\Given}{\,\middle|\,}

% (theorem environments removed — not used in current version)

\begin{document}

% --------------------
% Title
% --------------------
\title{
Contiguous-Segment Mixed-Precision Quantization for Asymmetric CSI Feedback Autoencoders
}

% --------------------
% Authors
% --------------------
\author{
\IEEEauthorblockN{Hyunjae Park and Wan Choi}
\IEEEauthorblockA{Department of Electrical and Computer Engineering, Seoul National University\\
Seoul 08826, South Korea\\
Email: \{hyunjae.park, wanchoi\}@snu.ac.kr}
}

\maketitle

\begin{abstract}
In frequency-division duplex (FDD) massive multiple-input multiple-output (MIMO) systems, downlink CSI must be estimated at the user equipment (UE) and fed back to the base station (BS). In practice, UE-side inference is constrained by latency, energy, and compute, making lightweight encoders and low-precision execution unavoidable. To handle long-tailed delay--angular CSI under strict UE budgets, we propose an asymmetric CSI feedback autoencoder that employs a lightweight state-space model (SSM) encoder at the UE, exploiting soft locality and natural quantization robustness, paired with a high-capacity Transformer decoder at the BS. Building on this architecture, we develop an interaction-aware segment-level mixed-precision quantization framework: a state-conditioned importance surface quantifies the segment-level degradation from contiguous-group perturbation of encoder blocks; in aggressive compression regimes where multiple blocks are simultaneously assigned low precision, segment-level perturbation captures inter-block interactions that per-block additive analysis may underestimate. Offline, segment-level dynamic programming constructs optimal contiguous segmentation policies under a fixed UE compute budget, and online deployment reduces to a single table lookup with no runtime optimization. A lightweight structural descriptor $\zeta$ conditions the degradation analysis on the wireless operating state, discriminating encoder-friendly and encoder-hostile CSI realizations with a 3.3\,dB NMSE spread. Experimental results demonstrate that the proposed asymmetric architecture improves the accuracy--complexity trade-off: at 93\% encoder BOPs saving, MT-AE achieves $-12.69$\,dB NMSE through the proposed Segment DP framework while CRNet collapses to $-6.79$\,dB under per-block ILP, a gap exceeding 5.9\,dB. Segment DP yields up to 0.8\,dB NMSE improvement over per-block integer linear programming under the same compute budget. An outage-aware extension, Joint DP, further optimises per-state budget allocation across channel conditions, shifting the outage cliff to higher savings at no additional NMSE cost.
\end{abstract}

\begin{IEEEkeywords}
CSI feedback, mixed-precision quantization, segment-level quantization, interaction-aware, massive MIMO, state-space model
\end{IEEEkeywords}

% =========================================================
\section{Introduction}
% =========================================================

\subsection{Background and Motivation}
In frequency-division duplex (FDD) massive multiple-input multiple-output (MIMO) systems, CSI feedback is essential for downlink precoding and rate adaptation. Unlike time-division duplex systems, FDD operation does not permit channel reciprocity; therefore, the UE must estimate downlink CSI and convey a compressed representation to the BS over a limited feedback link. Recent deep-learning-based CSI feedback methods employ autoencoders to learn compact latent representations that can be quantized and transmitted efficiently \cite{CsiNet,CsiNetPlus,CRNet,CLNet,TransNet,ACCsiNet}.

A practical bottleneck is UE-side feasibility at inference time. CSI compression must operate under strict latency and energy constraints, which often forces lightweight encoder architectures and low-precision execution. Under such operation, CSI distortion is unavoidable and its impact is not limited to reconstruction error: distorted CSI directly affects downlink transmission rates when used for precoding and link adaptation \cite{Jindal2006}. This motivates CSI feedback designs that jointly optimize the encoder architecture for quantization robustness and the mixed-precision policy for interaction-aware degradation control.

\subsection{Asymmetric CSI Feedback Architecture}
In the delay--angular domain, CSI often exhibits localized energy around dominant propagation paths while retaining non-negligible long-range components due to off-grid effects and finite DFT resolution. Encoders that impose hard locality (e.g., limited receptive fields) may effectively truncate such components under strict budgets, incurring irreversible structural loss. In contrast, state-space model (SSM) encoders aggregate inputs through soft memory with exponentially decaying influence, allowing long-range dependencies to be attenuated rather than abruptly truncated. Moreover, SSM encoders with contractive state dynamics provide natural robustness to low-precision execution. Recent work also highlights emerging SSM opportunities in wireless communications \cite{Mamba,MambaWireless}.

Motivated by this structural match and the inherent UE--BS computational asymmetry, we propose an asymmetric CSI feedback architecture that places a lightweight SSM-based encoder at the UE and a higher-capacity Transformer decoder at the BS.

\subsection{Interaction-Aware Mixed-Precision Quantization}
On top of the fixed asymmetric architecture, we develop an interaction-aware segment-level mixed-precision quantization framework, which constructs optimal precision assignments with minimal online complexity.

The framework defines a state-conditioned importance surface that measures the segment-level degradation from contiguous-group perturbation of encoder blocks, conditioned on a wireless state summarizing SNR and CSI structure. Per-block perturbation---where a single block is quantised while all others remain at anchor precision---provides a tractable additive approximation that is often sufficient at moderate compression levels. However, in aggressive compression regimes where multiple blocks are simultaneously assigned low precision, joint degradation effects may emerge that the additive model underestimates; segment-level perturbation captures these inter-block interactions. Offline, segment-level dynamic programming constructs optimal contiguous segmentation policies under a fixed compute budget. Online, the UE performs a single table lookup to deploy the pre-computed policy with no runtime optimization.

\subsection{Contributions}
The main contributions of this paper are:
\begin{itemize}
\item An asymmetric SSM--Transformer CSI feedback architecture that pairs a lightweight Mamba encoder at the UE with a Transformer decoder at the BS, providing strong quantisation robustness through contractive state dynamics.
\item \emph{Segment DP}: a segment-level interaction-aware mixed-precision formulation with dynamic programming; contiguous joint perturbation captures inter-block degradation that per-block additive analysis underestimates in aggressive compression regimes, and a dynamic program finds the optimal contiguous segmentation under a fixed compute budget, yielding up to 0.8\,dB improvement over per-block ILP.
\item \emph{Joint DP}: an outage-aware extension that optimises per-state budget allocation across channel conditions (indexed by the structural descriptor $\zeta$), improving outage probability at no additional NMSE cost by redistributing compute toward difficult channels.
\item A lightweight structural descriptor $\zeta$ that discriminates encoder-friendly and encoder-hostile CSI realizations (3.3\,dB NMSE spread), enabling state-conditioned degradation analysis with $\mathcal{O}(1)$ online complexity.
\end{itemize}

\noindent
The remainder of the paper is organized as follows. Section~\ref{sec:system_model} describes the system model. Section~\ref{sec:asym_arch} presents the asymmetric architecture. Sections~\ref{sec:offline} and~\ref{sec:online} develop the offline and online stages of the mixed-precision framework, respectively. Section~\ref{sec:experiments} reports experimental results, and Section~\ref{sec:conclusion} concludes.

% =========================================================
\section{System Model and Problem Formulation}
\label{sec:system_model}
% =========================================================

\subsection{System Model}
\label{subsec:system_model}
We consider an FDD massive MIMO downlink system where a BS with $N_t$ transmit antennas serves a UE with $N_r$ receive antennas. Let $\mathbf{H}[n] \in \mathbb{C}^{N_r \times N_t}$ denote the downlink MIMO channel on the $n$-th OFDM subcarrier, $n=0,\ldots,N_f-1$.

Applying an $N_f$-point unitary inverse DFT (IDFT) across the $N_f$ subcarriers yields the delay-domain channel taps
\begin{equation}
\mathbf{H}_d[\ell]
= \frac{1}{\sqrt{N_f}} \sum_{n=0}^{N_f-1} \mathbf{H}[n] e^{j 2\pi n\ell/N_f},
\quad \ell=0,\ldots,N_f-1.
\label{eq:delay_idft}
\end{equation}
Each tap is mapped to the angular domain using unitary DFT matrices
\begin{equation}
\mathbf{X}[\ell] = \mathbf{F}_r \mathbf{H}_d[\ell] \mathbf{F}_t^{\mathsf{H}},
\label{eq:ang_transform}
\end{equation}
where $\mathbf{F}_r \in \mathbb{C}^{N_r \times N_r}$ and $\mathbf{F}_t \in \mathbb{C}^{N_t \times N_t}$ are unitary DFT matrices.

Due to limited delay spread, most energy is concentrated in the first $N_a \le N_f$ delay taps. We retain
\begin{equation}
\mathbf{X}_a \triangleq \big[ \mathbf{X}[0], \mathbf{X}[1], \ldots, \mathbf{X}[N_a-1] \big].
\label{eq:truncate}
\end{equation}

At the UE, an encoder $f_\theta(\cdot)$ maps $\mathbf{X}_a$ to a latent vector
\begin{equation}
\mathbf{z} = f_\theta(\mathbf{X}_a), \quad \mathbf{z}\in\mathbb{R}^{D}, \quad D \ll 2 N_a N_r N_t.
\label{eq:encoder}
\end{equation}
(Complex CSI is represented in a real-valued format, e.g., by stacking real and imaginary parts; we omit this bookkeeping for clarity.)

The compression ratio is defined as
\begin{equation}
\mathrm{CR} \triangleq \frac{D}{2 N_a N_r N_t}.
\label{eq:cr}
\end{equation}

The latent is quantized and transmitted over the feedback link,
\begin{equation}
\tilde{\mathbf{z}} = Q_{\mathrm{fb}}(\mathbf{z}),
\label{eq:latent_quant}
\end{equation}
and the BS reconstructs the truncated CSI as
\begin{equation}
\hat{\mathbf{X}}_a = g_\phi(\tilde{\mathbf{z}}).
\label{eq:decoder}
\end{equation}

\subsection{UE Constraints and Design Perspective}
\label{subsec:constraints}
In practical deployments, CSI feedback is constrained by UE-side latency, energy, and compute limits, while the BS typically has substantially greater computational resources. This UE--BS asymmetry motivates architectures with a lightweight UE-side encoder and a higher-capacity BS-side decoder.

We study the joint design of:
\begin{enumerate}
\item a UE-feasible asymmetric encoder--decoder architecture that improves the accuracy--complexity trade-off; and
\item an interaction-aware segment-level mixed-precision quantization framework for the UE-side encoder that minimises distortion under a fixed UE compute budget.
\end{enumerate}
Section~\ref{sec:asym_arch} describes the architecture, and Sections~\ref{sec:offline}--\ref{sec:online} develop the mixed-precision scheduling framework.

% =========================================================
\section{Asymmetric CSI Feedback Architecture}
\label{sec:asym_arch}
% =========================================================

\subsection{Long-Tailed Locality of Delay--Angular Domain CSI}
\label{subsec:longtail}
In the delay--angular domain, CSI energy is often concentrated around dominant angular components associated with physical propagation paths. However, locality is not strictly compact: when a physical angle $\theta_0$ falls between DFT bins, the angular-domain coefficient at bin $k$ equals $|X[k]|=|a_0|\,|D_{N_t}(\theta_0-k/N_t)|$, where
\begin{equation}
D_N(f)=\frac{\sin(\pi N f)}{N\sin(\pi f)}
\label{eq:dirichlet_kernel}
\end{equation}
is the Dirichlet kernel. For offset $d\triangleq|k-k_0|\ge 1$, the standard bound $|\sin(\pi f)|\ge 2|f|$ for $|f|\le 1/2$ yields $|D_N|\le 1/(\pi N|f|)$, giving the oscillatory long tail:
\begin{equation}
x_{i_0+d} \;\approx\; c_0\,\frac{\eta_d}{1+d},
\qquad
d \triangleq |i-i_0|\ge 0,
\label{eq:dirichlet_osc_tail}
\end{equation}
where $\eta_d$ is an oscillatory factor satisfying $|\eta_d|\le 1$, and $c_0>0$ absorbs the path gain and normalization constants. In particular, the magnitude envelope obeys
\begin{equation}
|x_i|
\;\lesssim\;
\frac{c_0}{1+|i-i_0|}.
\label{eq:longtail_bound}
\end{equation}
Thus, distant angular coefficients can remain non-negligible, and strict hard locality may cause irreversible structural loss under tight UE budgets.

\subsection{Hard-Locality Encoding and Structural Truncation Loss}
\label{subsec:hard_locality}
Encoders based on CNNs impose finite receptive fields and thus a hard-locality constraint. Let $L$ denote an effective receptive-field radius. Under the envelope in \eqref{eq:longtail_bound}, the residual tail energy beyond radius $L$ satisfies
\begin{align}
E_{\mathrm{hard}}(L)
\triangleq
\sum_{|i-i_0|>L} |x_i|^2
\;\lesssim\;
c_0^2\sum_{d=L+1}^{\infty}\frac{1}{(1+d)^2}
\;\le\;
\frac{c_0^2}{1+L}
=
\mathcal{O}\!\left(L^{-1}\right),
\quad d \triangleq |i-i_0|.
\label{eq:hard_tail}
\end{align}
The bound indicates that reducing truncation loss requires substantially increasing $L$ (e.g., deeper/wider CNNs), which increases UE-side compute and latency.

\subsection{Soft-Locality Encoding via State-Space Aggregation}
\label{subsec:ssm_soft}
State-space model (SSM) encoders provide soft locality: rather than hard-truncating long-range coefficients, they recursively aggregate inputs with exponentially decaying aggregation weights. A linear time-invariant SSM can be written as
\begin{equation}
\mathbf{s}_{k+1}=\mathbf{A}\mathbf{s}_k+\mathbf{B}\mathbf{u}_k,
\quad
\mathbf{y}_k=\mathbf{C}\mathbf{s}_k,
\label{eq:ssm}
\end{equation}
which yields the convolutional form
\begin{equation}
\mathbf{y}_k=\sum_{d\ge 0}\mathbf{C}\mathbf{A}^{d}\mathbf{B}\mathbf{u}_{k-d}.
\label{eq:ssm_conv}
\end{equation}
If $\mathbf{A}$ is stable, then there exist constants $c_1>0$ and $\alpha>0$ such that $\|\mathbf{A}^{d}\|\le c_1 e^{-\alpha d}$ for all $d\ge 0$.
To compare hard truncation and soft aggregation on the same footing, define the aggregated tail contribution beyond horizon $L$:
\begin{equation}
R_{\mathrm{soft}}(L)
\triangleq
\left\|
\mathbf{C}\sum_{d>L}\mathbf{A}^{d}\mathbf{B}\mathbf{u}_{k-d}
\right\|.
\label{eq:soft_contrib_def}
\end{equation}
Motivated by Dirichlet-type leakage, we model the input envelope as $\|\mathbf{u}_{k-d}\|\lesssim \frac{c_0}{1+d}$ (up to oscillations as in \eqref{eq:dirichlet_osc_tail}). Then
\begin{align}
R_{\mathrm{soft}}(L)
&\le
\|\mathbf{C}\|\|\mathbf{B}\|\sum_{d>L}\|\mathbf{A}^{d}\|\,\|\mathbf{u}_{k-d}\|
\;\lesssim\;
c_3\sum_{d=L+1}^{\infty}\frac{e^{-\alpha d}}{1+d}
=
\mathcal{O}\!\left(\frac{e^{-\alpha L}}{1+L}\right),
\label{eq:soft_tail_bound}
\end{align}
for some constant $c_3>0$. This quantifies that the incremental influence of very distant coefficients is suppressed under exponential aggregation weights, while those coefficients are still retained and recursively aggregated rather than being abruptly truncated.

\subsection{Quantization Robustness}
\label{subsec:ssm_quant_robust}

UE-side inference requires low-precision execution, and post-training quantization introduces parameter perturbations $\hat{\boldsymbol{\theta}}=\boldsymbol{\theta}+\Delta\boldsymbol{\theta}$.
In Mamba's parameterization, the discretized state matrix is diagonal with entries $\bar{A}_{t,ii}=\exp(-\Delta_{t,i}\,\mathrm{softplus}(A_{\mathrm{param},i}))$. Since both $\Delta_{t,i}>0$ and $\mathrm{softplus}(\cdot)>0$, every entry satisfies $0<\bar{A}_{t,ii}<1$, making the state recursion inherently contractive for all inputs.
Unlike feedforward stacking where quantization errors accumulate with depth, a contractive state recursion prevents perturbation amplification: the output deviation converges to a finite steady-state bound independent of sequence length.
This robustness property is structurally consistent with the soft-locality analysis in Section~\ref{subsec:ssm_soft}: the same exponential decay enabling efficient long-tail modeling also stabilizes quantization-induced distortions under UE-side low-precision execution.

\subsection{Asymmetric Architecture}
\label{subsec:asym_conclude}

Based on these observations, we adopt an asymmetric CSI feedback architecture that deploys a lightweight Mamba-style SSM encoder at the UE \cite{Mamba, VMamba} and a higher-capacity Transformer-based decoder at the BS \cite{TransNet}.
The UE encoder benefits from SSM's linear-time complexity and quantization robustness, while the BS decoder leverages Transformer's expressive capacity without UE-side constraints.
This asymmetry is a practical design choice matched to the CSI feedback channel's inherent compute imbalance between UE and BS.

% =========================================================
\section{Interaction-Aware Mixed-Precision Quantization: Offline Policy Construction}
\label{sec:offline}
% =========================================================

This section develops the offline stage of the mixed-precision framework. The UE-side encoder contains $M$ quantizable blocks (contiguous parameter groups in the FC bottleneck layer). Mixed precision is applied to encoder weights while activations remain fixed at INT16. The feedback-link quantization $Q_{\mathrm{fb}}(\cdot)$ in \eqref{eq:latent_quant} is kept fixed. The framework constructs a \emph{state-conditioned importance surface} by measuring segment-level importance from contiguous-group perturbation, and derives optimal contiguous segmentation policies from this surface via dynamic programming.

\subsection{Mixed-Precision Policy Space}
\label{subsec:policy_space}

Let the UE-side encoder consist of $M$ quantizable blocks (contiguous parameter groups in the FC bottleneck layer). A \emph{policy} $\pi=\{(S_q,b_q)\}_{q=1}^{Q}$ assigns bit-widths $b_q\in\mathcal{B}$ to contiguous segments $S_q$ forming a partition of the $M$ quantizable FC chunks in the encoder bottleneck, with normalized cost
\begin{equation}
\kappa_\pi = \sum_{q=1}^{Q} \kappa(S_q, b_q) \in (0,1],
\label{eq:kappa_policy_sum}
\end{equation}
where $\kappa(S_q, b_q)$ denotes the compute cost of executing segment $S_q$ at bit-width $b_q$ (e.g., BOPs), normalized such that $\kappa_{\pi^{(\mathrm{FP32})}}=1$, and $\mathcal{B}$ denotes the discrete set of supported bit-widths.

As a special case, block-wise assignment $\pi=(b_1,\ldots,b_M)$ treats each block as a singleton segment.
The induced block-wise policy space has cardinality
\begin{equation}
|\Pi| = |\mathcal{B}|^M,
\label{eq:policy_cardinality}
\end{equation}
which makes exhaustive evaluation infeasible for moderate $M$.

\subsection{Wireless State}
\label{subsec:wireless_state}

We define a two-dimensional wireless state $\boldsymbol{\xi}=(\rho,\zeta)$ summarizing the current operating condition, where $\rho$ is the received SNR and $\zeta$ captures CSI structure.

For the structural component, define the normalized energy map
\begin{equation}
P_{ij}(\mathbf{X}_a) = \frac{|X_{ij}|^2}{\sum_{u,v}|X_{uv}|^2},
\label{eq:energy_map}
\end{equation}
and soft-locality kernels aligned with the encoder's receptive structure:
\begin{equation}
[K_d]_{i,i'} = \frac{(1+|i-i'|)^{-\alpha_d}}{\sum_{i''}(1+|i-i''|)^{-\alpha_d}}, \qquad
[K_a]_{j,j'} = \frac{(1+|j-j'|)^{-\alpha_a}}{\sum_{j''}(1+|j-j''|)^{-\alpha_a}},
\label{eq:kernels}
\end{equation}
where $\alpha_d,\alpha_a>0$ are decay exponents along the delay and angular axes.
The \emph{full} structural mismatch index is
\begin{equation}
\zeta_{\mathrm{full}}(\mathbf{X}_a) = 1 - \operatorname{tr}\!\left(P^\top K_d\, P\, K_a\right),
\label{eq:zeta_full}
\end{equation}
which requires two matrix products costing $\mathcal{O}(N_a^2 N_t + N_a N_t^2)$.

For efficient UE-side computation, we introduce a \emph{proxy} mismatch index based on per-axis marginals.
Define the delay and angular marginal energy distributions
\begin{equation}
[\mathbf{p}_d]_i = \sum_{j} P_{ij}, \qquad
[\mathbf{p}_a]_j = \sum_{i} P_{ij},
\label{eq:marginals}
\end{equation}
and per-axis compactness scores
\begin{equation}
c_d = \mathbf{p}_d^\top K_d\, \mathbf{p}_d, \qquad
c_a = \mathbf{p}_a^\top K_a\, \mathbf{p}_a.
\label{eq:per_axis_compactness}
\end{equation}
The proxy mismatch index is
\begin{equation}
\zeta(\mathbf{X}_a) \triangleq 1 - \lambda_d\, c_d - \lambda_a\, c_a,
\label{eq:struct_state}
\end{equation}
where $\lambda_d, \lambda_a > 0$ with $\lambda_d + \lambda_a = 1$.
The proxy requires only marginalization and two matrix--vector products, reducing cost to $\mathcal{O}(N_a N_t + N_a^2 + N_t^2)$.
Table~\ref{tab:zeta_runtime} summarizes the computational comparison.

Since $\zeta$ serves solely as a lookup key for the policy map $\mu$ in \eqref{eq:policy_map}, only the rank ordering of realizations needs to be preserved; the mapping is invariant to monotone transformations of $\zeta$.
We adopt equal weights $\lambda_d = \lambda_a = 1/2$ as the default.
Alternative choices include dimension-weighted ($\lambda_d = N_t/(N_a+N_t)$, $\lambda_a = N_a/(N_a+N_t)$) and OLS-fitted weights minimizing $\sum_i(\zeta_{\mathrm{full},i} - \zeta_{\mathrm{proxy},i})^2$ on the training set; these are compared in ablation.

\begin{table}[t]
\centering
\caption{Computational cost of the structural mismatch index ($N_a = N_t = 32$).}
\label{tab:zeta_runtime}
\setlength{\tabcolsep}{5pt}
\begin{tabular}{lcc}
\toprule
Variant & Complexity & MACs \\
\midrule
$\zeta_{\mathrm{full}}$ & $\mathcal{O}(N_a^2 N_t + N_a N_t^2)$ & 65{,}536 \\
$\zeta_{\mathrm{proxy}}$ & $\mathcal{O}(N_a N_t + N_a^2 + N_t^2)$ & 4{,}096 \\
\bottomrule
\end{tabular}
\end{table}

Small $\zeta$ indicates compact, encoder-friendly CSI; large $\zeta$ indicates diffuse multipath or leakage tails.

\textbf{Remark} (Data-driven weight optimization).
The constraint $\lambda_d + \lambda_a = 1$ defines a one-dimensional family.
Given profiling data, the weights can be optimized as
$(\lambda_d^*, \lambda_a^*) = \argmin_{\lambda_d + \lambda_a = 1} \sum_{i}(\zeta_{\mathrm{full},i} - \zeta_{\mathrm{proxy},i})^2$,
or more directly by maximizing Spearman rank correlation with $\zeta_{\mathrm{full}}$.
This data-driven extension is left for future investigation.

\subsection{Distortion Metric}
\label{subsec:shortfall}

Let $\pi_{\mathrm{anc}}$ denote a high-precision anchor policy.
The distortion metric used for importance measurement is the NMSE delta:

\textbf{NMSE delta.}
The NMSE increase under policy $\pi$ relative to the anchor is
\begin{equation}
\Delta L_i(\pi) = \mathrm{NMSE}_i(\pi) - \mathrm{NMSE}_i(\pi_{\mathrm{anc}}),
\label{eq:nmse_delta}
\end{equation}
which directly measures reconstruction degradation and serves as the optimisation objective throughout the framework.

\textbf{Remark} (alternative metrics).
We also investigated a rate-based reliability shortfall
\begin{equation}
V_i(\pi;\gamma) = \frac{[\gamma\, r_i^{\mathrm{ref}} - r_i(\pi)]_+}{\gamma\, r_i^{\mathrm{ref}} + \eta_r},
\label{eq:shortfall}
\end{equation}
where $r_i^{\mathrm{ref}}$ is the reference rate under the anchor and $\gamma\in(0,1)$ is a reliability target.
In our experiments, shortfall-based and NMSE-delta-based policies produced comparable results for the single-user system considered here; shortfall may offer additional benefits in multi-user scenarios where per-user rate guarantees are required.

We write $\Delta\ell_i(\pi)$ generically for the chosen distortion increase, and the importance framework proceeds identically regardless of metric choice.

\subsection{State-Conditioned Importance Surface}
\label{subsec:importance_surface}

\textbf{Why segment-level perturbation.}
Per-block perturbation---where a single encoder block is quantised while all others remain at anchor precision---provides a tractable additive approximation of quantisation sensitivity. At moderate compression levels, this approximation is often sufficient. However, in aggressive compression regimes where multiple blocks are simultaneously assigned low precision, joint degradation effects may emerge that the additive model underestimates. In particular, the nonlinear sensitivity of beamforming gain to directional CSI errors only manifests when multiple blocks are simultaneously perturbed. To capture such interactions, we adopt segment-level perturbation.

\textbf{Segment-level perturbation.}
To capture inter-block interactions, we perturb contiguous \emph{segments} $S=[l{:}r)$ of consecutive blocks simultaneously.
Let $\pi^{(S\to b)}$ denote the policy that assigns all blocks in segment $S$ to bit-width $b$ while keeping the remaining blocks at anchor precision.
The segment-level distortion increase is
\begin{equation}
\Delta\ell_i(S,b) = \Delta\ell_i(\pi^{(S\to b)}),
\label{eq:segment_delta}
\end{equation}
where $\Delta\ell_i$ is the chosen distortion metric (NMSE delta or shortfall).
Because multiple blocks are perturbed simultaneously, the resulting error enters the nonlinear regime in which directional CSI errors disproportionately degrade beamforming gain---an effect that single-block perturbation may underestimate.

The state-conditioned segment importance surface is defined as
\begin{equation}
\widehat{\Omega}_S(\boldsymbol{\xi}; b) \triangleq \E\!\left[
\Delta\ell_i(S,b) \Given \boldsymbol{\xi}_i = \boldsymbol{\xi}
\right],
\label{eq:segment_omega}
\end{equation}
measuring the expected distortion increase from assigning segment $S$ to bit-width $b$, conditioned on wireless state $\boldsymbol{\xi}=(\rho,\zeta)$.

This is the core quantity of the framework. Because $\widehat{\Omega}_S(\boldsymbol{\xi}; b)$ is conditioned on both SNR and CSI structure, it automatically encodes regime-dependent behavior: segments that are NMSE-sensitive under high SNR and segments whose degradation is amplified under diffuse structure are distinguished by the surface without manual regime separation.

\subsection{Importance Estimation}
\label{subsec:estimation}

The state space $(\rho,\zeta)$ is discretized into a $J\times K$ grid of quantile cells $\{C_{j,k}\}$, with $J$ bins along $\rho$ and $K$ bins along $\zeta$.
For each candidate segment $S=[l{:}r)$ with $|S|\le L_{\max}$ and each bit-width $b\in\mathcal{B}$, the segment importance is estimated from the profiling data.

\textbf{Raw estimator.}
Within each cell $(j,k)$, the raw segment importance is the cell-wise average:
\begin{equation}
\widehat{\Omega}_{S,b}^{\mathrm{raw},(j,k)} = \frac{1}{N_{jk}} \sum_{i:\,\boldsymbol{\xi}_i\in C_{j,k}} \Delta\ell_i(S,b).
\label{eq:raw_omega}
\end{equation}

\textbf{Shrinkage stabilization.}
For cells with limited samples, we shrink toward the global segment-bit mean:
\begin{equation}
\widehat{\Omega}_{S,b}^{(j,k)} = \alpha_{jk}\,\widehat{\Omega}_{S,b}^{\mathrm{raw},(j,k)} + (1-\alpha_{jk})\,\bar{\Omega}_{S,b},
\label{eq:shrinkage}
\end{equation}
where $\bar{\Omega}_{S,b} = \frac{1}{N}\sum_{i=1}^{N}\Delta\ell_i(S,b)$ is the global mean and $\alpha_{jk}=N_{jk}/(N_{jk}+\tau_{\mathrm{shr}})$ controls shrinkage intensity.

\textbf{Monotone calibration.}
A 2-D isotonic calibration is applied to enforce physical monotonicity:
\begin{itemize}
\item \emph{Nonincreasing in $\rho$}: higher SNR reduces quantization sensitivity.
\item \emph{Nondecreasing in $\zeta$}: diffuse CSI structure increases quantization sensitivity.
\end{itemize}
At runtime, bilinear interpolation or direct cell lookup yields $\widehat{\Omega}_{S,b}(\boldsymbol{\xi}_t)$.

\subsection{Policy Optimization via Dynamic Programming}
\label{subsec:policy_opt}

Given the segment importance surface $\widehat{\Omega}_S(\boldsymbol{\xi}; b)$ constructed in Section~\ref{subsec:importance_surface} and a fixed compute budget~$c$, the optimal contiguous segmentation is found by dynamic programming:
\begin{equation}
F(m,c) = \min_{\substack{0\le l < m \\ b\in\mathcal{B}}}
\Big[F(l,\;c{-}\kappa_{[l:m)}^{(b)}) + \widehat{\Omega}_{[l:m)}(\boldsymbol{\xi};b)\Big],
\label{eq:dp}
\end{equation}
with base case $F(0,c)=0$ for all $c\ge 0$, where $\kappa_{[l:m)}^{(b)}$ is the normalised BOPs cost of segment $[l{:}m)$ at precision~$b$.
The DP runs in $\mathcal{O}(M^2 |\mathcal{B}| C)$ time, where $C$ is the budget discretisation granularity.
Unlike per-block ILP, this formulation naturally respects the contiguous structure of the FC weight matrix and captures the interaction effects measured by $\widehat{\Omega}_S$.

The offline stage produces a state-to-policy lookup table
\begin{equation}
\mu(\boldsymbol{\xi}, \bar{c}) = \pi^*(\boldsymbol{\xi}),
\label{eq:policy_map}
\end{equation}
mapping each wireless state and budget level to the optimal contiguous segmentation and per-segment precision assignment.

\textbf{Joint DP: outage-aware budget allocation.}
The Segment DP framework admits a natural outer budget allocation that we term \emph{Joint DP}. The inner Segment DP produces a distortion--budget curve $D_j(c) = F_j(M,c)$ for each state cell~$j$. An outer problem distributes a system-wide average budget~$\bar{c}$ across state cells:
\begin{equation}
\min_{\{c_j\}}\;\sum_{j=1}^{J} p_j\, D_j(c_j)
\qquad\text{s.t.}\quad
\sum_{j=1}^{J} p_j\, c_j \le \bar{c},
\label{eq:outer}
\end{equation}
where $p_j$ is the fraction of samples in bin~$j$.
This water-filling-like allocation redistributes compute toward difficult channel conditions (high $\zeta$), improving outage probability at no additional NMSE cost. Section~\ref{subsec:offline_ablation} validates this extension experimentally.

% =========================================================
\section{Online Deployment}
\label{sec:online}
% =========================================================

Given the offline-constructed policy map $\mu$ in \eqref{eq:policy_map}, the UE performs a single table lookup per CSI realization.

\subsection{Online Decision Rule}
\label{subsec:online_rule}

At runtime, the UE computes the current wireless state $\boldsymbol{\xi}_t=(\rho_t,\zeta(\mathbf{X}_t))$, looks up the pre-computed policy $\pi_t=\mu(\boldsymbol{\xi}_t)$, and executes the encoder under the selected mixed-precision policy.
No online optimization is performed: the entire adaptation is reduced to a single table lookup.
Algorithm~\ref{alg:online_deploy} summarizes the procedure.

\begin{algorithm}[t]
\caption{Online Policy Deployment}
\label{alg:online_deploy}
\begin{algorithmic}[1]
\REQUIRE CSI $\mathbf{X}_a$, SNR $\rho_t$, offline policy map $\mu$
\STATE Compute $\zeta_t = \zeta(\mathbf{X}_a)$ via \eqref{eq:struct_state}
\STATE Form wireless state $\boldsymbol{\xi}_t = (\rho_t,\, \zeta_t)$
\STATE Look up policy $\pi_t \leftarrow \mu(\boldsymbol{\xi}_t)$
\STATE Execute encoder $\tilde{\mathbf{z}}_t = f_\theta(\mathbf{X}_a;\, \pi_t)$
\RETURN quantized feedback $\tilde{\mathbf{z}}_t$
\end{algorithmic}
\end{algorithm}

\subsection{Online Complexity}
\label{subsec:online_complexity}

The online cost consists of computing $\zeta(\mathbf{X}_t)$ via \eqref{eq:struct_state} ($\mathcal{O}(N_a N_t + N_a^2 + N_t^2)$ for the proxy; see Table~\ref{tab:zeta_runtime}) and a table lookup ($\mathcal{O}(1)$), both negligible relative to encoder inference. Therefore, the framework enables interaction-aware mixed-precision adaptation with no measurable overhead.

% =========================================================
\section{Experimental Results}
\label{sec:experiments}
% =========================================================
This section validates the proposed asymmetric CSI feedback framework along three dimensions: (i) structural efficiency of the asymmetric UE encoder, (ii) effectiveness of segment-level dynamic programming versus per-block ILP, and (iii) the $\zeta$ descriptor's ability to discriminate degradation regimes under identical fixed UE-side compute budgets.

\subsection{Experimental Setup}
\label{subsec:exp_setup}
We consider an FDD massive MIMO downlink system with $N_t=32$ transmit antennas at the BS and a single-antenna UE ($N_r=1$). CSI is represented in the delay--angular domain with $N_f=1024$ OFDM subcarriers, and only the first $N_a=32$ delay taps are retained for feedback. The BS employs maximum ratio transmission (MRT) based on reconstructed CSI, and performance is evaluated at SNR levels of 10, 20, and 30~dB.

We use COST~2100 outdoor channel realizations \cite{COST2100} (carrier frequency $f_c=300$~MHz) and additionally evaluate generalization on FR1 TDL-A/B/C profiles from 3GPP TR~38.901 \cite{3GPP38901}. The UE encoder is a Mamba-style selective SSM with 2-D cross-scan \cite{Mamba, VMamba}, and the BS decoder is Transformer-based. Training uses AdamW with learning rate $10^{-3}$ for 200 epochs. Mixed-precision quantization is applied to UE-side encoder weights with candidate bit-widths $\{16,8,4,2\}$, while internal activations use 16-bit and the latent uses 8-bit quantization for feedback. Unless stated otherwise, complexity is reported as encoder FLOPs (full precision) or encoder BOPs (quantized inference), normalized consistently across methods. We report UE-side compute budgets as \emph{encoder BOP saving relative to an FP32 baseline} (32-bit weights and 32-bit activations); specifically, a saving level $S\in\mathcal{S}$ corresponds to an average cost constraint $\bar{c}=1-S$, and the optimization problems are formulated over the corresponding cost set $\mathcal{C}=\{1-S: S\in\mathcal{S}\}$. Table~\ref{tab:sim_setup} summarizes all simulation parameters.

\begin{table}[t]
\centering
\caption{Simulation parameters.}
\label{tab:sim_setup}
\setlength{\tabcolsep}{5pt}
\begin{tabular}{ll}
\toprule
\textbf{Parameter} & \textbf{Value} \\
\midrule
System & FDD downlink, $N_t=32$, $N_r=1$ \\
CSI representation & Delay--angular, $N_f=1024$, $N_a=32$ \\
SNR (dB) & 10 / 20 / 30 \\
Channel model & COST~2100, outdoor, $f_c=300$~MHz \cite{COST2100} \\
TDL evaluation & FR1 TDL-A/B/C \cite{3GPP38901} \\
\midrule
UE encoder & Mamba SSM \cite{Mamba, VMamba} \\
BS decoder & Transformer \\
Compression ratio & $\mathrm{CR}=1/4$ \\
Optimizer & AdamW, lr $10^{-3}$, 200 epochs, batch 1000 \\
\midrule
Weight precision & 16 / 8 / 4 / 2 bit (mixed) \\
Activation precision & 16-bit (internal), 8-bit (latent) \\
Saving range $\mathcal{S}$ & 85--97\% vs.\ FP32, step 0.05\% \\
Outage threshold $\gamma$ & 0.99 / 0.98 / 0.95 \\
\bottomrule
\end{tabular}
\end{table}

\subsection{Structural Efficiency of the Asymmetric Encoder}
\label{subsec:struct_eff}
We first evaluate structural reconstruction performance without encoder weight quantization. Table~\ref{tab:struct_cr14_final} reports NMSE and computational complexity at $\mathrm{CR}=1/4$.

\begin{table}[t]
\centering
\caption{FP32 baseline reconstruction and complexity ($\mathrm{CR}=1/4$, outdoor).}
\label{tab:struct_cr14_final}
\setlength{\tabcolsep}{4pt}
\begin{tabular}{lccc}
\toprule
Model & NMSE (dB) & Enc.\ FLOPs (M) & Total FLOPs (M) \\
\midrule
CsiNet   & -8.75  &  1.09 &   5.41 \\
CsiNet+  & -12.40 &  1.45 &  24.57 \\
CRNet    & -12.71 &  1.20 &   5.12 \\
CLNet    & -12.87 &  1.35 &   4.05 \\
TransNet & -14.86 & 17.83 &  35.72 \\
\textbf{MT-AE (Ours)} & \textbf{-15.37} & \textbf{4.70} & 22.53 \\
\bottomrule
\end{tabular}
\end{table}


\subsection{Quantization Robustness under Uniform Precision}
\label{subsec:quant_robust}
We evaluate uniform weight quantization under INT16/INT8/INT4 at $\mathrm{CR}=1/4$.
Table~\ref{tab:uniform_quant_cr4} reports NMSE and encoder BOPs for each precision level.
Fig.~\ref{fig:uniform_rpmpq}(a) visualizes the accuracy--efficiency trade-off: MT-AE retains near-lossless NMSE down to INT8 ($87.5\%$ BOPs saving), whereas the CNN baselines degrade sharply beyond INT16.
Fig.~\ref{fig:uniform_rpmpq}(b) shows the offline mixed-precision Pareto frontier, where MT-AE consistently achieves the best NMSE--efficiency trade-off across the full saving range.
Fig.~\ref{fig:segment_dp_all_models} extends this comparison to all architectures under the proposed Segment DP framework, confirming that MT-AE maintains its advantage across the full saving range.
This robustness gap motivates the mixed-precision strategy developed in Sections~\ref{subsec:offline_ablation} and \ref{subsec:online_outage}.

\begin{table}[t]
\centering
\caption{Uniform weight quantization ($\mathrm{CR}=1/4$, outdoor). Activations are fixed at INT16; listed precision refers to weight bit-width. $^\dagger$Quantization collapse (NMSE $\ge 0$\,dB).}
\label{tab:uniform_quant_cr4}
\setlength{\tabcolsep}{4pt}
\begin{tabular}{llccc}
\toprule
Model & Weight & NMSE (dB) & Enc.\ BOPs (M) & Saving (\%) \\
\midrule
CsiNet & INT16 & $-8.74$  & 268.61 & 75.00 \\
       & INT8  &   1.46$^\dagger$   & 134.30 & 87.50 \\
       & INT4  &  19.40$^\dagger$   &  67.15 & 93.75 \\
\midrule
CRNet  & INT16 & $-12.71$ & 268.61 & 75.00 \\
       & INT8  &  $-3.57$ & 134.31 & 87.50 \\
       & INT4  &   10.36$^\dagger$  &  67.15 & 93.75 \\
\midrule
CLNet  & INT16 & $-12.82$ & 268.69 & 75.00 \\
       & INT8  &   0.15$^\dagger$   & 134.34 & 87.50 \\
       & INT4  &  23.36$^\dagger$   &  67.17 & 93.75 \\
\midrule
\textbf{MT-AE (Ours)} & INT16 & $\mathbf{-15.37}$ & 1245.46 & 75.00 \\
             & INT8  & $\mathbf{-15.19}$ &  622.85 & 87.50 \\
             & INT4  &   0.03$^\dagger$   &  311.47 & 93.75 \\
\bottomrule
\end{tabular}
\end{table}

\begin{figure*}[t]
\centering
\includegraphics[width=0.95\textwidth]{figures/fig_uniform_rpmpq.pdf}
\caption{NMSE under varying BOPs saving relative to FP32 ($\mathrm{CR}=1/4$, outdoor).
(a)~Uniform weight quantization (activations fixed at INT16):
MT-AE maintains near-lossless NMSE up to $87.5\%$ saving (INT8), whereas CNN baselines collapse beyond INT16.
(b)~Offline mixed-precision Pareto frontier ($85$--$97\%$ saving): MT-AE achieves the best NMSE--efficiency trade-off across the full range.}
\label{fig:uniform_rpmpq}
\end{figure*}

\begin{figure}[t]
\centering
\includegraphics[width=0.95\columnwidth]{figures/fig_segment_dp_all_models.pdf}
\caption{Segment DP Pareto frontier across all architectures ($\mathrm{CR}=1/4$, outdoor). MT-AE maintains the best NMSE--efficiency trade-off; at 93\% saving, MT-AE achieves $-12.69$\,dB while CRNet degrades to $-6.79$\,dB.}
\label{fig:segment_dp_all_models}
\end{figure}


\subsection{Segment DP vs.\ HAWQ-ILP and Joint DP Extension}
\label{subsec:offline_ablation}

Fig.~\ref{fig:kl_vs_ilp} presents a two-panel ablation comparing the proposed DP-based methods against baselines on MT-AE.

\textbf{Panel~(a): NMSE improvement via Segment DP.}
Panel~(a) compares the per-layer HAWQ-ILP baseline with the proposed Segment DP.
HAWQ-ILP treats each layer independently, whereas Segment DP captures inter-block interactions via contiguous-segment dynamic programming under the BOPs budget.
Segment DP consistently achieves lower NMSE, with the gap widening beyond $90\%$ saving where multiple blocks are simultaneously assigned low precision.
Joint DP, which redistributes compute across $\zeta$ bins, achieves nearly identical NMSE to Segment DP, confirming that the outage-aware allocation incurs no additional distortion cost.

\textbf{Panel~(b): Outage improvement via Joint DP.}
Panel~(b) extends to the outage dimension ($\gamma=0.99$, SNR${}=20$\,dB): we compare the Segment DP policy with equal budget allocation across all $\zeta$ bins against Joint DP, which additionally optimises the per-state budget allocation to minimise population-weighted outage under an average compute constraint.
Joint DP shifts the outage cliff to higher savings by redistributing compute budget toward difficult channel conditions (high $\zeta$), demonstrating that outage-aware allocation provides meaningful improvement at no additional NMSE cost.
Table~\ref{tab:budget_validation} confirms budget consistency.

\begin{figure}[t]
\centering
\includegraphics[width=0.95\columnwidth]{figures/kl_vs_ilp.pdf}
\caption{Ablation on MT-AE ($\mathrm{CR}=1/4$, outdoor).
(a)~NMSE: Segment DP (proposed) vs.\ per-layer HAWQ-ILP and Joint DP; Segment DP consistently achieves the lowest NMSE, with the gap over HAWQ-ILP widening beyond $90\%$ saving. Joint DP overlaps Segment DP, confirming no additional NMSE cost.
(b)~Outage ($\gamma=0.99$, SNR${}=20$\,dB): Joint DP (outage-aware per-state budget allocation) vs.\ Segment DP with equal allocation across $\zeta$ bins; Joint DP shifts the outage cliff from $87\%$ to $88\%$ saving by redistributing compute toward difficult channels.}
\label{fig:kl_vs_ilp}
\end{figure}

\subsection{Outage Performance under Identical Budgets}
\label{subsec:online_outage}
We compare uniform precision, static mixed precision (offline-only), and the proposed Segment DP framework under identical fixed UE-side encoder BOP budgets. Performance is evaluated via outage probability under targets $\gamma \in \{0.99, 0.98, 0.95\}$ at SNR levels of 10, 20, and 30~dB. To highlight the cross-architecture advantage of our framework, we also compare against per-block ILP applied to CNN-based encoders (CRNet, CLNet, CsiNet) under identical saving levels.

\subsubsection{Transmission Rate and Outage Definition}
To compute per-subcarrier MRT precoding, the reconstructed delay--angular CSI $\hat{\mathbf{X}}_a$ is mapped back to the frequency domain by inverse angular transforms and a DFT across delay taps. Specifically, we form $\hat{\mathbf{H}}_d[\ell]=\mathbf{F}_r^{\mathsf{H}}\hat{\mathbf{X}}[\ell]\mathbf{F}_t$ for $\ell=0,\ldots,N_a-1$, zero-pad $\hat{\mathbf{H}}_d[\ell]=\mathbf{0}$ for $\ell\ge N_a$, and compute $\hat{\mathbf{H}}[n]=\frac{1}{\sqrt{N_f}}\sum_{\ell=0}^{N_f-1}\hat{\mathbf{H}}_d[\ell]e^{-j2\pi n\ell/N_f}$. For $N_r=1$, the corresponding channel vector is $\hat{\mathbf{h}}[n]\in\mathbb{C}^{N_t}$. For the $t$-th CSI realization, the achievable rate is defined as
\begin{equation}
r_t(\mathbf{x})
\triangleq
\frac{1}{N_f}\sum_{n=0}^{N_f-1}
\log_2\!\left(1+\rho \abs{\mathbf{h}_t[n]^{\mathsf{H}} \mathbf{w}_{\mathbf{x}}[n]}^2 \right),
\label{eq:rate_def}
\end{equation}
where $\mathbf{h}_t[n]$ is the true channel, $\rho$ is the received SNR, and $\mathbf{w}_{\mathbf{x}}[n] \triangleq \mathbf{x}[n]/\norm{\mathbf{x}[n]}_2$ is the MRT beamforming vector. Equal power allocation across subcarriers is assumed.

Under policy $\pi$, let $\hat{\mathbf{h}}_{t,\pi}[n]$ denote the reconstructed channel used to form the MRT precoder, and define $r_t(\pi) \triangleq r_t(\hat{\mathbf{h}}_{t,\pi})$. The reference rate is $r_t^{\mathrm{ref}} \triangleq r_t(\mathbf{h}_t^{(a)})$, where $\mathbf{h}_t^{(a)}[n]$ is obtained from the truncated ground-truth $\mathbf{X}_a$ via the same inverse transforms (with zero-padding), so that $r_t^{\mathrm{ref}}$ does not penalize the representation truncation itself. An outage event is declared when
\begin{equation}
r_t(\pi) < \gamma\, r_t^{\mathrm{ref}}.
\label{eq:outage_event}
\end{equation}

\subsubsection{Outage Curves}
Fig.~\ref{fig:outage_final} presents rate-based outage probability versus BOPs saving budget for three SNR levels and three reliability thresholds ($\gamma \in \{0.99, 0.98, 0.95\}$). The proposed Segment DP (solid) consistently shifts the outage cliff to higher savings compared to per-block ILP (dashed), indicating that the interaction-aware formulation provides improved distortion control at aggressive saving levels.
This improvement is consistent with the SSM encoder's inherent quantization robustness: the contractive state recursion limits perturbation amplification, and the segment-level formulation captures the inter-block interactions that emerge when multiple blocks operate at low precision simultaneously.

\subsubsection{Cross-Architecture Comparison}
The advantage of the proposed framework is amplified when compared across architectures (Fig.~\ref{fig:cross_arch_outage}). At 93\% encoder BOPs saving, MT-AE with Segment DP achieves $-12.69$\,dB NMSE, whereas CRNet under per-block ILP collapses to $-6.79$\,dB---a gap exceeding 5.9\,dB. This gap arises from two compounding factors: (i)~the SSM encoder's contractive state dynamics provide inherent robustness to aggressive quantization, and (ii)~Segment DP captures inter-block interactions that per-block ILP misses. CNN-based encoders lack the contractive recursion and suffer disproportionate degradation when multiple FC chunks are simultaneously assigned low precision.

\begin{figure*}[t]
\centering
\includegraphics[width=0.95\textwidth]{figures/fig_cross_arch_outage.pdf}
\caption{Cross-architecture outage comparison under the proposed Segment DP framework ($\mathrm{CR}=1/4$, outdoor). Three panels correspond to reliability thresholds $\gamma \in \{0.99, 0.98, 0.95\}$. MT-AE maintains low outage probability at aggressive saving levels where CNN-based encoders collapse.}
\label{fig:cross_arch_outage}
\end{figure*}

\begin{figure*}[t]
\centering
\includegraphics[width=0.95\textwidth]{figures/fig_online_outage.pdf}
\caption{Rate-based outage probability under varying BOPs saving relative to FP32 ($\mathrm{CR}=1/4$, outdoor). Each panel corresponds to a different downlink SNR. Solid: Segment DP; dashed: per-block ILP. Colors indicate the reliability threshold $\gamma$.}
\label{fig:outage_final}
\end{figure*}

\subsection{Budget Consistency of Joint DP}
\label{subsec:budget_validation}
Table~\ref{tab:budget_validation} validates that the Joint DP outer budget allocation enforces the long-term compute constraint: the achieved average saving deviates by at most $0.022\%$ from the target, confirming that the water-filling-like allocation is practically exact.

\begin{table}[t]
\centering
\caption{Budget consistency of Joint DP per-state budget allocation (SNR\,=\,20\,dB, $\gamma=0.99$).}
\label{tab:budget_validation}
\setlength{\tabcolsep}{6pt}
\begin{tabular}{ccc}
\toprule
Target Saving (\%) & Achieved Saving (\%) & Deviation (\%) \\
\midrule
$87.5\%$ & $87.48\%$ & $0.022$ \\
$90.0\%$ & $90.00\%$ & $0.000$ \\
$92.5\%$ & $92.49\%$ & $0.007$ \\
\bottomrule
\end{tabular}
\end{table}

\subsection{Validation under FR1 TDL Profiles and UE Runtime}
\label{subsec:tdl_runtime}
To evaluate generalization to standardized propagation conditions, we train a single MT-AE on mixed FR1 TDL-A/B/C realizations and evaluate per profile on a UE platform. Table~\ref{tab:tdl_generalization} reports NMSE under FP32, uniform INT4, and offline MPQ, together with measured encoder latency and BOPs saving. We also report per-sample UE-side runtime in Table~\ref{tab:runtime_validation}; the offline lookup table construction is a one-time step performed before deployment.

\begin{table}[t]
\centering
\caption{MT-AE performance on FR1 TDL profiles ($\mathrm{CR}=1/4$).
Single model trained on mixed TDL-A/B/C data; evaluated per profile on a UE platform. (To be filled.)}
\label{tab:tdl_generalization}
\setlength{\tabcolsep}{4pt}
\begin{tabular}{lccccc}
\toprule
Profile & NMSE & NMSE & NMSE & Latency & BOPs \\
        & (FP32) & (Unif.\ INT4) & (Offline MPQ) & (ms) & Saving \\
\midrule
TDL-A & \textemdash & \textemdash & \textemdash & \textemdash & \textemdash \\
TDL-B & \textemdash & \textemdash & \textemdash & \textemdash & \textemdash \\
TDL-C & \textemdash & \textemdash & \textemdash & \textemdash & \textemdash \\
\bottomrule
\end{tabular}
\end{table}

\begin{table}[t]
\centering
\caption{Per-sample UE-side latency: FP32 baseline vs.\ Segment DP. Offline lookup table construction is performed once before deployment. (To be filled.)}
\label{tab:runtime_validation}
\setlength{\tabcolsep}{4pt}
\begin{tabular}{lcc}
\toprule
Component & FP32 & Segment DP \\
\midrule
Encoder inference  & \textemdash\,ms & \textemdash\,ms \\
Policy selection   & ---             & \textemdash\,$\mu$s \\
Total per sample   & \textemdash\,ms & \textemdash\,ms \\
Avg.\ BOPs saving  & $0\%$           & \textemdash\,\% \\
\bottomrule
\end{tabular}
\end{table}

% =========================================================
\section{Conclusion}
\label{sec:conclusion}
% =========================================================
This paper studied CSI feedback for FDD massive MIMO under UE-side inference constraints.
An asymmetric autoencoder pairing a lightweight SSM encoder at the UE with a Transformer decoder at the BS provides a quantization-robust baseline.
The first methodological contribution, Segment DP, is a contiguous-segment interaction-aware dynamic program that finds optimal mixed-precision segmentations under a fixed compute budget: by perturbing contiguous groups of encoder blocks jointly, the formulation captures inter-block degradation that per-block additive analysis underestimates in aggressive compression regimes, yielding up to 0.8\,dB NMSE improvement over per-block ILP.
The second contribution, Joint DP, extends Segment DP by optimising per-state budget allocation across channel conditions, shifting the outage cliff to higher savings at no additional NMSE cost.
Cross-architecture evaluation confirms the compounding advantage: at 93\% encoder BOPs saving, MT-AE with Segment DP achieves $-12.69$\,dB NMSE while CRNet collapses to $-6.79$\,dB under per-block ILP, a gap exceeding 5.9\,dB attributable to both the SSM encoder's contractive robustness and the segment-level formulation's ability to capture inter-block interactions. A lightweight structural descriptor $\zeta$ discriminates encoder-friendly and encoder-hostile CSI realizations with a 3.3\,dB NMSE spread, enabling $\mathcal{O}(1)$ online policy lookup.

The current framework applies post-training quantization with fixed activation precision (INT16) and uses equal proxy weights ($\lambda_d=\lambda_a=1/2$). Several directions are left for future work: (i) rate-aware distortion objectives (e.g., reliability shortfall), which produced comparable results to NMSE delta in our single-user experiments but may offer additional benefits in multi-user scenarios; (ii) extensions to multi-antenna UEs, joint weight--activation optimization, data-driven proxy weight optimization, and non-stationary channel statistics.
% --------------------
% References
% --------------------
\begin{thebibliography}{99}

\bibitem{CsiNet}
C.-K. Wen, W.-T. Shih, and S.~Jin,
``Deep learning for massive MIMO CSI feedback,''
\emph{IEEE Wireless Communications Letters}, vol.~7, no.~5, pp.~748--751, Oct.~2018.

\bibitem{CRNet}
J.~Guo, C.-K. Wen, S.~Jin, and G.~Y.~Li,
``Multi-resolution CSI feedback with deep learning in massive MIMO system,''
\emph{IEEE Journal on Selected Areas in Communications}, vol.~40, no.~8, pp.~2499--2514, Aug.~2022.

\bibitem{CLNet}
S.~Ji and M.~Li,
``CLNet: Complex input lightweight neural network designed for massive MIMO CSI feedback,''
\emph{IEEE Wireless Communications Letters}, vol.~10, no.~11, pp.~2318--2322, Nov.~2021.

\bibitem{TransNet}
Z.~Liu, Y.~Chen, and C.-K. Wen,
``TransNet: Full-stack deep learning-based CSI feedback for massive MIMO systems,''
\emph{IEEE Transactions on Wireless Communications}, vol.~21, no.~10, pp.~8261--8275, Oct.~2022.

\bibitem{CsiNetPlus}
J.~Guo, C.-K. Wen, S.~Jin, and G.~Y.~Li,
``Convolutional neural network-based multiple-rate compressive sensing for massive MIMO CSI feedback: Design, simulation, and analysis,''
\emph{IEEE Transactions on Wireless Communications}, vol.~19, no.~4, pp.~2827--2840, Apr.~2020.

\bibitem{ACCsiNet}
B.~Cao, Y.~Yang, P.~Ran, D.~He, and G.~He,
``ACCsiNet: Asymmetric convolution-based autoencoder framework for massive MIMO CSI feedback,''
\emph{IEEE Communications Letters}, vol.~25, no.~12, pp.~3873--3877, Dec.~2021.

\bibitem{Jindal2006}
N.~Jindal,
``MIMO broadcast channels with finite-rate feedback,''
\emph{IEEE Transactions on Information Theory}, vol.~52, no.~11, pp.~5045--5060, Nov.~2006.


\bibitem{COST2100}
L.~Liu \emph{et al.},
``The COST 2100 MIMO channel model,''
\emph{IEEE Wireless Communications}, vol.~19, no.~6, pp.~92--99, Dec.~2012.

\bibitem{Mamba}
A.~Gu and T.~Dao,
``Mamba: Linear-time sequence modeling with selective state spaces,''
arXiv preprint arXiv:2312.00752, 2023.

\bibitem{VMamba}
Y.~Liu, Y.~Tian, Y.~Zhao, H.~Yu, L.~Xie, Y.~Wang, Q.~Ye, and Y.~Liu,
``VMamba: Visual state space model,''
arXiv preprint arXiv:2401.10166, 2024.

\bibitem{MambaWireless}
R.~Zhang, R.~Zhang, Y.~Lu, W.~Chen, B.~Ai, and D.~Niyato,
``Mamba for wireless communications and networking: Principles and opportunities,''
arXiv preprint arXiv:2508.00403, Aug.~2025, doi: 10.48550/arXiv.2508.00403.

\bibitem{3GPP38901}
3GPP,
``Study on channel model for frequencies from 0.5 to 100~GHz,''
3GPP TR~38.901, V16.1.0, Dec.~2020.

\end{thebibliography}

\end{document}
